\exercisesheader{}

% 3

\eoce{\qt{Polusi udara dan luaran kelahiran, komponen penelitian\label{study_components_airpoll}}
Para peneliti mengumpulkan data untuk mengkaji hubungan antara polutan udara
dan kelahiran prematur di California Selatan. Selama penelitian, tingkat polusi
udara diukur oleh stasiun pemantau kualitas udara. Secara khusus, kadar karbon
monoksida dicatat dalam bagian per juta, nitrogen dioksida dan ozon dalam bagian
per seratus juta, dan partikulat kasar (PM$_{10}$) dalam $\mu g/m^3$.
Data lama kehamilan dikumpulkan untuk 143,196 kelahiran antara tahun 1989
dan 1993, lalu paparan polusi udara selama kehamilan dihitung untuk setiap
kelahiran. Analisis menunjukkan bahwa peningkatan PM$_{10}$ di udara sekitar dan,
pada tingkat yang lebih rendah, konsentrasi CO mungkin berkaitan dengan kejadian
kelahiran prematur.\footfullcite{Ritz+Yu+Chapa+Fruin:2000}
\begin{parts}
\item Tentukan pertanyaan penelitian utama dalam penelitian ini.
\item Siapa subjek dalam penelitian ini, dan berapa banyak yang disertakan?
\item Apa saja variabel dalam penelitian ini? Tentukan apakah setiap variabel
bersifat numerik atau kategoris. Jika numerik, nyatakan apakah variabel tersebut
diskret atau kontinu. Jika kategoris, nyatakan apakah variabel tersebut ordinal.
\end{parts}
}{}

% 4

\eoce{\qt{Metode Buteyko, komponen penelitian\label{study_components_buteyko}}
Metode Buteyko adalah teknik pernapasan dangkal yang dikembangkan oleh Konstantin
Buteyko, seorang dokter Rusia, pada tahun 1952. Bukti anekdotal menunjukkan bahwa
metode Buteyko dapat mengurangi gejala asma dan meningkatkan kualitas hidup. Dalam
sebuah penelitian ilmiah untuk menentukan keefektifan metode ini, para peneliti
merekrut 600 pasien asma berusia 18-69 tahun yang bergantung pada obat untuk
penanganan asma. Pasien-pasien ini secara acak dibagi menjadi dua kelompok
penelitian: satu kelompok mempraktikkan metode Buteyko dan kelompok lainnya tidak.
Pasien diberi skor dari 0 hingga 10 untuk kualitas hidup, aktivitas, gejala asma,
dan pengurangan penggunaan obat. Rata-rata, peserta dalam kelompok Buteyko
mengalami penurunan gejala asma yang signifikan dan peningkatan kualitas hidup.\footfullcite{McDowan:2003}
\begin{parts}
\item Tentukan pertanyaan penelitian utama dalam penelitian ini.
\item Siapa subjek dalam penelitian ini, dan berapa banyak yang disertakan?
\item Apa saja variabel dalam penelitian ini? Tentukan apakah setiap variabel
bersifat numerik atau kategoris. Jika numerik, nyatakan apakah variabel tersebut
diskret atau kontinu. Jika kategoris, nyatakan apakah variabel tersebut ordinal.
\end{parts}
}{}

% 5

\eoce{\qt{Kecurangan, komponen penelitian\label{study_components_cheaters}}
Para peneliti yang mengkaji hubungan antara kejujuran, usia, dan pengendalian diri
melakukan eksperimen terhadap 160 anak berusia antara 5 dan 15 tahun.
Peserta melaporkan usia, jenis kelamin, dan apakah mereka anak tunggal.
Para peneliti meminta setiap anak untuk melempar koin seimbang secara pribadi,
mencatat hasilnya (putih atau hitam) pada selembar kertas, dan memberi tahu bahwa
hanya anak yang melaporkan putih yang akan mendapat hadiah.
Temuan penelitian dapat dirangkum sebagai berikut:
``Separuh anak secara tegas diberi tahu agar tidak berbuat curang, sedangkan
yang lain tidak diberi petunjuk tegas apa pun.
Dalam kelompok tanpa petunjuk, probabilitas berbuat curang ditemukan seragam
di antara kelompok yang dibentuk berdasarkan karakteristik anak.
Dalam kelompok yang secara tegas diberi tahu agar tidak berbuat curang,
anak perempuan lebih kecil kemungkinannya berbuat curang; sementara tingkat
kecurangan tidak bervariasi menurut usia pada anak laki-laki, tingkat tersebut
menurun seiring bertambahnya usia pada anak perempuan.''\footfullcite{Bucciol:2011}
\begin{parts}
\item Tentukan pertanyaan penelitian utama dalam penelitian ini.
\item Siapa subjek dalam penelitian ini, dan berapa banyak yang disertakan?
\item
  Berapa banyak variabel yang dicatat untuk setiap subjek dalam penelitian ini
  agar temuan tersebut dapat disimpulkan? Nyatakan variabel-variabel itu beserta
  jenisnya.
\end{parts}
}{}

\D{\newpage}

% 6

\eoce{\qt{Perilaku mengambil permen, komponen penelitian\label{study_components_stealers}}
Dalam sebuah penelitian tentang hubungan antara kelas sosioekonomi dan perilaku
tidak etis, 129 mahasiswa sarjana University of California, Berkeley diminta
mengidentifikasi diri sebagai anggota kelas sosial bawah atau atas dengan
membandingkan diri mereka dengan orang lain yang memiliki uang paling banyak
(paling sedikit), pendidikan paling tinggi (paling rendah), dan pekerjaan paling
dihormati (paling tidak dihormati). Mereka juga diberi sebuah stoples berisi
permen yang dibungkus satu per satu dan diberi tahu bahwa permen tersebut
diperuntukkan bagi anak-anak di laboratorium terdekat, tetapi mereka boleh
mengambil beberapa jika mau. Setelah menyelesaikan beberapa tugas yang tidak
berkaitan, peserta melaporkan jumlah permen yang mereka ambil.\footfullcite{Piff:2012}
\begin{parts}
\item Tentukan pertanyaan penelitian utama dalam penelitian ini.
\item Siapa subjek dalam penelitian ini, dan berapa banyak yang disertakan?
\item Penelitian menemukan bahwa mahasiswa yang diidentifikasi sebagai kelas atas
mengambil lebih banyak permen daripada mahasiswa lainnya. Berapa banyak variabel
yang dicatat untuk setiap subjek dalam penelitian ini agar temuan tersebut dapat
disimpulkan? Nyatakan variabel-variabel itu beserta jenisnya.
\end{parts}
}{}

% 7

\eoce{\qt{Migrain dan akupunktur,
    Bagian II\label{migraine_and_acupuncture_exp_resp}}
Latihan~\ref{migraine_and_acupuncture_intro}
memperkenalkan sebuah penelitian yang mengkaji apakah akupunktur
memiliki pengaruh terhadap migrain.
Para peneliti melakukan penelitian terkontrol acak
dengan menempatkan pasien secara acak ke salah satu dari dua kelompok:
perlakuan atau kontrol.
Pasien dalam kelompok perlakuan menerima akupunktur
yang dirancang secara khusus untuk menangani migrain.
Pasien dalam kelompok kontrol menerima akupunktur plasebo
(jarum ditusukkan pada lokasi yang bukan titik akupunktur).
Pasien ditanya 24 jam setelah menerima akupunktur
apakah mereka sudah bebas nyeri.
Apa variabel penjelas dan variabel respons dalam penelitian ini?
}{}

% 8

\eoce{\qt{Sinusitis dan antibiotik,
    Bagian II\label{sinusitis_and_antibiotics_exp_resp}}
Latihan~\ref{sinusitis_and_antibiotics_intro}
memperkenalkan sebuah penelitian yang mengkaji pengaruh pengobatan antibiotik
terhadap sinusitis akut.
Peserta penelitian menerima rangkaian pengobatan antibiotik selama 10 hari
(perlakuan) atau plasebo yang tampak dan terasa serupa (kontrol).
Pada akhir periode 10 hari, pasien ditanya apakah
gejala mereka membaik.
Apa variabel penjelas dan variabel respons dalam penelitian ini?
}{}

% 9

\eoce{\qt{Bunga iris Fisher\label{fisher_irises}} Sir Ronald Aylmer Fisher adalah
seorang ahli statistika, biologi evolusi, dan genetika asal Inggris yang bekerja
dengan sebuah himpunan data berisi panjang dan lebar sepal serta panjang dan lebar
petal dari tiga spesies bunga iris (\textit{setosa}, \textit{versicolor}, dan
\textit{virginica}). Himpunan data tersebut memuat 50 bunga dari setiap spesies.
\footfullcite{Fisher:1936} \\
\noindent\begin{minipage}[c]{0.48\textwidth}
\begin{parts}
\item Berapa banyak kasus yang disertakan dalam data?
\item Berapa banyak variabel numerik yang disertakan dalam data? Sebutkan
variabel-variabel tersebut dan tentukan apakah masing-masing kontinu atau diskret.
\item Berapa banyak variabel kategoris yang disertakan dalam data, dan apa saja
variabel tersebut? Daftarkan level (kategori) yang bersesuaian.
\end{parts}
\end{minipage}
\begin{minipage}[c]{0.01\textwidth}
\ 
\end{minipage}
\begin{minipage}[c]{0.2\textwidth}
\begin{center}
\Figures[Foto bunga iris ungu.]{}{eoce/fisher_irises}{irisversicolor}
\end{center}
\end{minipage}
\begin{minipage}[c]{0.01\textwidth}
\ 
\end{minipage}
\begin{minipage}[c]{0.23\textwidth}
{\raggedright\footnotesize Foto oleh Ryan Claussen
(\oiRedirect{textbook-flickr_ryan_claussen_iris_picture}{http://flic.kr/p/6QTcuX}) 
\href{https://creativecommons.org/licenses/by-sa/2.0/}{lisensi CC~BY-SA~2.0}}
\end{minipage}
}{}

% 10

\eoce{\qt{Kebiasaan merokok penduduk Britania Raya\label{smoking_habits_UK_datamatrix}} Sebuah survei
dilakukan untuk mempelajari kebiasaan merokok penduduk Britania Raya. Matriks
data di bawah menampilkan sebagian data yang dikumpulkan dalam survei ini.
Perhatikan bahwa ``$\pounds$'' berarti pound sterling Britania, ``cig'' berarti
rokok, dan ``N/A'' menunjukkan komponen data yang hilang. \footfullcite{data:smoking}
\begin{center}
\scriptsize{
\begin{tabular}{rccccccc}
\hline
	& kelamin 	 & usia 	& status 	& pendapatan 					     & merokok & akhirPekan	& hariKerja \\
\hline
1 	& Female & 42 	& Single 	& Under $\pounds$2,600 			     & Yes 	 & 12 cig/day   & 12 cig/day \\ 
2 	& Male	 & 44	& Single 	& $\pounds$10,400 to $\pounds$15,600 & No	 & N/A 			& N/A \\ 
3 	& Male 	 & 53 	& Married   & Above $\pounds$36,400 		     & Yes 	 & 6 cig/day 	& 6 cig/day \\ 
\vdots & \vdots & \vdots & \vdots & \vdots 				             & \vdots & \vdots 	    & \vdots \\ 
1691 & Male  & 40   & Single 	& $\pounds$2,600 to $\pounds$5,200   & Yes 	 & 8 cig/day 	& 8 cig/day \\   
\hline
\end{tabular}
}
\end{center}
\begin{parts}
\item Apa yang diwakili oleh setiap baris matriks data?
\item Berapa banyak peserta yang disertakan dalam survei?
\item Tentukan apakah setiap variabel dalam penelitian bersifat numerik atau
kategoris. Jika numerik, tentukan apakah variabel itu kontinu atau diskret.
Jika kategoris, nyatakan apakah variabel tersebut ordinal.
\end{parts}
}{}

% 11

\eoce{\qt{Bandara di AS\label{US Airports}}
Visualisasi di bawah menunjukkan sebaran geografis bandara di Amerika Serikat
daratan utama dan Washington, DC.
Visualisasi ini dibuat berdasarkan sebuah himpunan data dengan
setiap amatan berupa satu bandara.\footfullcite{data:usairports}
Batas kartografis pada peta bersumber dari U.S. Census Bureau.
\par\noindent\makebox[\textwidth][c]{%
\Figures[Empat salinan peta Amerika Serikat ditampilkan dalam kisi 2 $\times$ 2. Pada setiap peta, sumbu diberi label bujur (130 derajat barat hingga 60 derajat barat) dan lintang (20 derajat utara hingga 50 derajat utara). Kolom pertama grafik diberi label "penggunaan privat" dan kolom kedua "penggunaan publik". Baris pertama grafik diberi label "milik privat" dan baris kedua "milik publik". Keempat grafik menampilkan titik-titik yang masing-masing mewakili satu bandara. Ribuan titik tampak pada peta kiri atas (penggunaan privat, milik privat) dan peta kanan bawah (penggunaan publik, milik publik), sedangkan dua peta lainnya menampilkan relatif lebih sedikit titik, meskipun jumlahnya masih berkisar dari ratusan hingga beberapa ribu. Pada semua grafik, titik lebih padat di bagian tengah dan timur Amerika Serikat, lebih jarang di wilayah pegunungan dan gurun, lalu kembali lebih padat di sekitar negara bagian yang berbatasan dengan Samudra Pasifik, terutama di dekat kota-kota besar.]{0.85}{eoce/airports}{airports}
}\par
\begin{parts}
\item
    Daftarkan variabel yang digunakan untuk membuat visualisasi ini.
\item
    Tentukan apakah setiap variabel dalam penelitian bersifat numerik
    atau kategoris.
    Jika numerik, tentukan apakah variabel itu kontinu atau diskret.
    Jika kategoris, nyatakan apakah variabel tersebut ordinal.
\end{parts}
}{}

% 12

\eoce{\qt{Pemungutan suara PBB\label{unvotes}}
Visualisasi berikut menunjukkan pola suara Amerika Serikat, Kanada, dan Meksiko
mengenai beragam isu di Majelis Umum Perserikatan Bangsa-Bangsa. Untuk setiap
tahun 1946--2015, visualisasi ini menampilkan persentase pemungutan suara
tercatat ketika setiap negara memilih ya pada setiap isu. Data terdiri atas
amatan pasangan negara--tahun.\footfullcite{data:unvotes}
\par\noindent\makebox[\textwidth][c]{%
\Figures[Ditampilkan sebuah kisi diagram pencar dengan garis tren yang ditumpangkan untuk masing-masing dari tiga kelompok titik (berwarna hijau, biru, dan merah) pada setiap grafik. Kisi grafik terdiri atas 2 baris dan 3 kolom; grafik dalam deskripsi ini disebut dengan nomor 1 sampai 3 pada baris pertama dan 4 sampai 6 pada baris kedua. Pada semua grafik, sumbu horizontal menyatakan "tahun" (sekitar 1945 hingga sekitar 2018), sedangkan sumbu vertikal menyatakan "persentase ya" dengan nilai dari 0\% hingga 100\%. Setiap grafik merangkum pola pemungutan suara untuk satu topik berbeda di Majelis Umum PBB dan untuk negara Kanada (biru), Meksiko (hijau), serta Amerika Serikat (merah). Setiap grafik memuat titik dan garis tren lentur (nonlinier) yang disesuaikan dengan titik-titik tersebut. Dalam semua kasus kecuali Grafik 2 untuk "Kolonialisme", titik (data) relatif jarang pada 1940 hingga 1960 dibandingkan tahun-tahun berikutnya. Grafik 1 mewakili "Pengendalian senjata dan perlucutan senjata"; semua negara bermula rendah, antara 0\% dan 25\%, lalu naik cepat menjelang 1960 menjadi antara 25\% dan 95\%. AS tetap paling rendah (berkisar sekitar 25\% hingga 40\%), Kanada sedikit lebih tinggi pada 50\% hingga 70\%, dan Meksiko paling tinggi, biasanya antara 85\% dan 100\%. Grafik 2 berlabel "Kolonialisme"; garis tren bermula antara 50\% dan 80\%, kemudian AS turun mendekati 0\% pada 1980, Kanada berfluktuasi antara 25\% dan 60\% sepanjang periode, dan Meksiko naik mendekati 100\% pada 1980. Grafik 3 mewakili "Pembangunan ekonomi"; ketiga negara bermula sekitar 25\% hingga 40\%. AS turun menjadi sekitar 5\% pada 1990 sebelum naik menjadi 20\%; Kanada turun menjadi sekitar 25\% pada 1985, naik menjadi 50\% pada 2000, lalu turun lagi menjadi 25\%; dan Meksiko naik menjadi sekitar 100\% pada 1980 sebelum turun menjadi sekitar 85\%. Grafik 4 mewakili "Hak asasi manusia"; semua negara mengelompok di sekitar 65\% pada 1945, lalu AS turun menjadi 25\% pada 1975 dan berfluktuasi antara 10\% dan 30\% selama sisa periode, Kanada perlahan turun menjadi sekitar 15\%, dan Meksiko naik mendekati 100\% pada 1985 lalu perlahan turun menjadi sekitar 80\%. Grafik 5 mewakili "Senjata dan bahan nuklir"; semua negara bermula mendekati 0\% pada 1945, lalu AS sedikit naik tetapi umumnya berfluktuasi antara 15\% dan 40\%, Kanada naik menjadi sekitar 60\% pada 1965 sebelum turun dan berfluktuasi di sekitar 40\% hingga 50\%, dan Meksiko naik cepat menjadi sekitar 90\% pada 1970 lalu mendekati 100\% seiring waktu. Grafik 6 mewakili "Konflik Palestina"; semua negara bermula antara 50\% dan 75\%, lalu AS terus turun menjadi sekitar 10\% pada 1985 dan setelah itu mendekati 5\%, Kanada sedikit turun menjadi sekitar 35\% pada 1970 sebelum naik menjadi sekitar 70\% pada 2000 dan kemudian turun cepat mendekati 0\%, sedangkan Meksiko naik bertahap menjadi sekitar 95\% pada 1985 lalu bertahan kurang lebih tetap.]{0.85}{eoce/unvotes}{unvotes}
}\par
\begin{parts}
\item Daftarkan variabel yang digunakan untuk membuat visualisasi ini.
\item Tentukan apakah setiap variabel dalam penelitian bersifat numerik atau
kategoris. Jika numerik, tentukan apakah variabel itu kontinu atau diskret.
Jika kategoris, nyatakan apakah variabel tersebut ordinal.
\end{parts}
}{}
